% =====================================================================
% Soft-wall -> sharp-wall mapping of the SEM conductivity map.
% Drop-in text for PoreExplorer.tex (draft 2026-08-28):
%   (1) replaces the second paragraph of Results 3.1 ("Two comparisons
%       validate the map. ...") from "Second, against the sharp-walled
%       analytical reference" to the end of that paragraph;
%   (2) replaces Appendix C (\label{app:sigma}) in full;
%   (3) lists the other sentences of the draft that the result changes.
% Every number is reproduced by
%   python -m sem.scripts.soft_wall_mapping all --out validation/soft_wall_mapping
% (see validation/soft_wall_mapping/README.md). New labels: eq:deficit_t,
% eq:rim, eq:kappa, eq:semfull, fig:si_softwall. New references to add
% to the .bib (verify before submission): Norris & Sheng 1989 (flanged-pipe
% end correction 0.8216a), Rayleigh's aperture conductance 2 sigma a.
% =====================================================================

% ---------------------------------------------------------------------
% (1) Results 3.1, second paragraph, from "Second, ..."
% ---------------------------------------------------------------------
Second, against the sharp-walled analytical reference, the \emph{uncorrected} map conductance lies $8$--$20\%$ low, with the deviation largest at small diameter---the signature of a fixed inward wall displacement. Mapping the ramp onto its electrically equivalent sharp wall removes this deficit, but a quantitative comparison also has to correct the reference itself. Equation~\eqref{eq:kowalczyk} uses Hall's access resistance $\rho/2d$, which is exact for an aperture in a thin plate and $4.6\%$ too small for any membrane thicker than about a pore radius, where the exact end correction is that of a flanged tube (Appendix~\ref{app:sigma}). Evaluated at the area-conserving offset $(d-0.54~\text{nm},\,L+0.54~\text{nm})$, Equation~\eqref{eq:kowalczyk} therefore deviates from the converged map conductance by $0.0\%$ at $d = 5$~nm and $+2.0\%$ at $d = 25$~nm ($L = 20$~nm), and by up to $3.1\%$ across $d = 5$--$50$~nm and $L = 1$--$40$~nm. This drift is the share of the access term growing with $d$, not a failure of the soft-wall correction; the near-perfect agreement at $d = 5$~nm is itself a cancellation of three effects of opposite sign. With the exact finite-thickness access term and the equivalent sharp geometry of Appendix~\ref{app:sigma}---the equal-area pore diameter, the membrane faces displaced by the same deficit length, and a rounded-rim correction that follows from the corner singularity---the reference reproduces the converged map conductance to $0.015\%$ for $L \ge 20$~nm ($0.002\%$ for $d \ge 5$~nm) and $0.12\%$ over $d = 2$--$50$~nm and $L = 1$--$40$~nm, with no parameter fitted to the solid-state pores being benchmarked [Figure~\ref{fig:si_softwall}]. In the finite golden-aspect box of Section~\ref{sec:domain} the map conductance sits $-0.24\%$ to $+0.54\%$ from its infinite-reservoir value, so the residual against the reference in Figure~\ref{fig:analytical_validation}(A) is box truncation, not geometry. \PENDING{regenerate Figure~\ref{fig:analytical_validation}(A) against Equation~\eqref{eq:semfull}; the $\pm 2.5\%$ band becomes a $\pm 0.5\%$ band set by the box.}

% ---------------------------------------------------------------------
% (2) Appendix C, full replacement
% ---------------------------------------------------------------------
\section{Mapping the conductivity map onto a sharp-walled pore}\label{app:sigma}
This appendix derives the sharp-walled geometry that is electrically equivalent to the ramped conductivity map of a cylindrical solid-state pore, and the sharp-walled reference against which it is compared. The statements are checked against an axisymmetric finite-element solution of $\nabla\!\cdot\!(\sigma\nabla\phi) = 0$ with unbounded reservoirs, Richardson-extrapolated to $\sim\!10^{-6}$ in conductance and confirmed by an independent finite-volume solver (\texttt{sem.axisym\_conductance}), and against an exact mode-matching solution of the sharp cylinder (\texttt{sem.cylinder\_mode\_matching}).

\paragraph{Conductivity ramp.}
The conductivity is $\sigma(\mathbf{r}) = \sigma_0\,f(D(\mathbf{r}))$, where $D$ is the Euclidean distance to the membrane and $f$ ramps linearly from $0$ at $r_\mathrm{min}$ to $1$ at $r_\mathrm{cut}$:
\begin{equation}
f(D) =
\begin{cases}
0, & D \le r_\mathrm{min},\\[2pt]
\dfrac{D - r_\mathrm{min}}{r_\mathrm{cut} - r_\mathrm{min}}, & r_\mathrm{min} < D < r_\mathrm{cut},\\[8pt]
1, & D \ge r_\mathrm{cut},
\end{cases}
\qquad
\begin{aligned}
r_\mathrm{min} &= 1.3~\text{\AA},\\
r_\mathrm{cut} &= 4.1~\text{\AA}.
\end{aligned}
\label{eq:framp}
\end{equation}
On a flat wall the ramp removes the conductance of a layer of thickness
\begin{equation}
t = \int_0^\infty \bigl[1 - f(D)\bigr]\,\mathrm{d}D = \tfrac{1}{2}(r_\mathrm{min}+r_\mathrm{cut}) = 2.7~\text{\AA},
\label{eq:deficit_t}
\end{equation}
the natural candidate for the wall displacement. The resistance of the pore has three parts---the bore, the membrane faces with the access region they bound, and the rim where they meet---and the ramp enters each differently.

\paragraph{Bore: equal-area radius, exact.}
Away from the pore ends, where the end modes have decayed as $e^{-3.83\,|z - z_\mathrm{end}|/a}$, $\sigma$ depends only on $r$ and $\phi = -Ez$ satisfies both $\nabla\!\cdot\!(\sigma\nabla\phi) = 0$ and the insulating condition on the wall. The field is axial, the conducting layers add in parallel, and the channel conductance per unit length is $\sigma_0\int_0^a f\,2\pi R\,\mathrm{d}R = \sigma_0\pi a_e^{2}$ with
\begin{equation}
a_e^{2} = (a - t)^{2} + \frac{\Delta r^{2}}{12},\qquad \Delta r = r_\mathrm{cut} - r_\mathrm{min},
\label{eq:matching}
\end{equation}
valid for $a \ge r_\mathrm{cut}$. Writing $f$ as a superposition of sharp steps, $t$ is the mean and $\Delta r^{2}/12 = 0.653$~\AA$^{2}$ the variance of the step position. No current crosses the iso-$\sigma$ surfaces, so no normal (``series'') contribution enters. The variance moves $a - a_e$ from $2.700$ to $2.685$~\AA\ at $d = 5$~nm, $+0.12\%$ in conductance.

\paragraph{Faces and access region: the same offset.}
Outside the pore the current converges radially onto the mouth (the ``$\mathrm{d}r$'' part of the access resistance), and along the membrane faces it runs tangentially, parallel to the ramp. A boundary-layer expansion shows that a flat wall carrying any ramp is equivalent to a sharp wall at $D = t$ with no correction at second order in the ramp width. The first normal-direction term is third order, $\partial_n\phi = -(\Delta r^{3}/48)\,\partial_x^{4}\phi$ at $n = t$, and is negligible. For a sharp-cornered membrane, moving both faces out by $t$ is exactly a change of thickness, $L \to L + 2t$. The radial convergence of the access current therefore needs no separate treatment: it is accounted for by the same deficit length.

\paragraph{Rim: a rounded corner.}
Because $D$ is a Euclidean distance, the iso-$\sigma$ contours wrap the rim as quarter circles, so the equivalent wall is the rounded (Minkowski) offset of the membrane rather than a sharp-cornered one. The sharp rim is a $270^{\circ}$ fluid wedge, whose potential is singular as $\phi \simeq K\varrho^{2/3}\cos(2\theta/3)$ at distance $\varrho$ from the edge, with $K = k\,E\,a^{1/3}$ and $E = I/\sigma_0\pi a^{2}$. Matching the ramp's local corner solution to this singular field shortens each end of the equivalent pore by
\begin{equation}
\Delta\ell = C\left[\frac{k(\lambda)}{k_\infty}\right]^{2} t^{4/3}a_e^{-1/3},\qquad C = -2\pi\beta\,k_\infty^{2},\qquad \lambda = \frac{L + 2t}{a_e}.
\label{eq:rim}
\end{equation}
The exponent $4/3 = 2\cdot\tfrac{2}{3}$ is that of the corner singularity, and $\beta$ is a constant of the local, two-dimensional corner problem: $\beta = -0.1494$ for the linear ramp and $-0.1272$ for a sharp wall on the $D = t$ iso-surface. $k(\lambda)$ is the rim strength of the sharp cylinder, $k_\infty = 0.7549$ for $L \gtrsim a$, growing as $0.5607\,\lambda^{-1/6}$ in thin membranes, where the plate edge begins to look like a knife edge. Together they give $C = 0.535$ for the default ramp, confirmed to $10^{-4}$ by a thick-membrane solve, and $0.455$ for the step. Neither $\beta$ nor $k$ is fitted to the conductance solutions benchmarked below. $\Delta\ell$ is $0.71$~\AA\ per end at $d = 5$~nm and $0.40$~\AA\ at $d = 25$~nm ($L = 20$~nm), so the thickness offset is $2(t - \Delta\ell) = 0.40$--$0.46$~nm rather than $0.54$~nm. For other linear ramps $C$ depends on the shape only through $\Delta r/t$, from $0.73$ for $r_\mathrm{min} = 0$ to $0.455$ for a step.

\paragraph{The sharp-walled reference: Hall's access term is a thin-membrane limit.}
Writing the resistance of a sharp-walled cylinder of radius $a$ as
\begin{equation}
R_\mathrm{sharp} = \frac{\rho}{\pi a^{2}}\bigl[L + 2\,\kappa(L/a)\,a\bigr],
\label{eq:kappa}
\end{equation}
Equation~\eqref{eq:kowalczyk} corresponds to $\kappa = \pi/4 = 0.7854$, the end correction of an equipotential disc. This is exact for an aperture in a thin plate ($L \to 0$), where the aperture is equipotential by symmetry. In a membrane of finite thickness it is not. The exact $\kappa$, computed by Ritz mode matching (a variational upper bound converged to $10^{-9}$), rises as $\kappa = \pi/4 + (\lambda/8)[\ln(2\pi/\lambda) - 3] + O(\lambda^{2})$ and reaches the flanged-tube value $\kappa_\infty = 0.82167$ (cf.\ $0.8216$, Norris and Sheng \PENDING{cite}) already at $L \approx a$. It is within $5\times10^{-4}$ of $\kappa_\infty$ for $L \ge a$, and approaches it as $e^{-j_{1,1}L/a}$, the decay of the first non-trivial mode of an insulating cylinder ($j_{1,1} = 3.8317$). Hall's access term is therefore $4.6\%$ too small for practically every solid-state pore, and Equation~\eqref{eq:kowalczyk} overestimates the conductance by up to $3.2\%$ (at $L \approx 0.2\,d$), $2.3\%$ at $L = 0.8\,d$, and $0.8\%$ at $L = 4d$. A closed form that is exact at both limits and within $5\times10^{-5}$ in between is
\begin{equation}
\kappa(\lambda) = \kappa_\infty - \bigl(\kappa_\infty - \tfrac{\pi}{4}\bigr)\,\tfrac{1}{2}\,e^{-j_{1,1}\lambda}\bigl[1 + e^{-4.313\,\lambda^{2/3}}\bigr].
\end{equation}
For $L \ge d/2$ it suffices to replace $1/d$ in Equation~\eqref{eq:kowalczyk} by $1.046/d$.

\paragraph{Combined mapping.}
The open-pore resistance of the conductivity map of a cylinder of nominal diameter $d = 2a$ and thickness $L$ is then
\begin{equation}
R_\mathrm{SEM} = \frac{\rho}{\pi a_e^{2}}\bigl[L + 2t + 2\,\kappa(\lambda)\,a_e - 2\Delta\ell\bigr],
\label{eq:semfull}
\end{equation}
with $a_e$ from Equation~\eqref{eq:matching}, $\lambda = (L + 2t)/a_e$, and $\Delta\ell$ from Equation~\eqref{eq:rim}. Against the converged map conductance it is accurate to $0.015\%$ for $L \ge 20$~nm ($0.002\%$ for $d \ge 5$~nm), $0.05\%$ for $L \ge 5$~nm, and $0.12\%$ over $d = 2$--$50$~nm and $L = 1$--$40$~nm. The largest residuals are in the thinnest membranes, where $L + 2t$ is only a few times $t$ and the corner expansion is strained. For ramps with $r_\mathrm{cut} = 3$--$6$~\AA\ and $r_\mathrm{min} = 0$--$2.5$~\AA\ the same expression holds to $0.03\%$ for $d \ge 3$~nm. Each ingredient matters [Figure~\ref{fig:si_softwall}(B,C)]. Hall's access term at the area-conserving offset errs by up to $3.1\%$, and by $-10.6\%$ at $d = 1.5$~nm, $L = 1$~nm. The exact access term without the rim correction errs by up to $-7.4\%$ for thin membranes and small pores. At $d = 5$~nm, $L = 20$~nm, the draft's near-perfect agreement decomposes into $-0.67\%$ (Hall's access term), $+0.12\%$ (the variance in Equation~\eqref{eq:matching}) and $+0.59\%$ (the rim). At $d = 25$~nm the access term alone accounts for $+2.2\%$, partly offset by the rim ($-0.2\%$).

The offset needed to force Equation~\eqref{eq:kowalczyk} onto the map conductance therefore grows with $d$ because of the access term. In the continuum it reaches $0.77$~nm at $d = 25$~nm when applied to both $d$ and $L$, or $0.86$~nm when applied to $d$ alone. The golden-aspect box adds $\approx 0.04$~nm, which accounts for most of the $0.94$~nm reported in Figure~\ref{fig:si_sigma}. The diameter offset of the mapping itself stays at $2(a - a_e) = 0.537$--$0.540$~nm, set by the ramp alone.

\paragraph{Reconciliation and robustness.}
The same area-conservation relation links the parametric map to an all-atom membrane: inverting Equation~\eqref{eq:matching} slice by slice, $a(z) = t + \sqrt{r_\mathrm{eq}(z)^{2} - \Delta r^{2}/12}$ with $r_\mathrm{eq}(z)=\sqrt{A_\mathrm{eff}(z)/\pi}$ measured on the all-atom map, gives the parametric radius that reproduces the all-atom open cross-section; for the drilled membranes of Table~\ref{tab:open_pore_current_match} it lies $\approx 0.9$~\AA\ inside the atom-centre drill radius. That comparison is between two soft walls built from the same ramp, so no offset enters; the analytical comparison needs the full mapping because Equation~\eqref{eq:kappa} assumes a sharp wall. At the reference point $d = 10$~nm, $L = 20$~nm, $\sigma = 10.5$~S/m, Equation~\eqref{eq:semfull} gives $26.18$~nS. For comparison, the nominal sharp pore conducts $29.61$~nS, and Equation~\eqref{eq:kowalczyk} at the area-conserving offset gives $26.39$~nS. The golden-aspect three-dimensional value, $26.3$~nS, lies $\approx 0.5\%$ above Equation~\eqref{eq:semfull}, consistent with the box truncation measured above. The inverse map---the nominal SEM geometry that conducts like a given sharp-walled pore, for instance an experimental one---follows from Equation~\eqref{eq:semfull} by fixed-point iteration (\texttt{sem.open\_pore\_theory.sem\_cylinder\_for\_sharp\_pore}).

\begin{figure}[!t]
\centering
\includegraphics[width=0.98\columnwidth]{SI/soft_wall_mapping.pdf}
\caption{\textbf{Mapping the conductivity map onto sharp-walled theory.} Converged axisymmetric solutions with unbounded reservoirs. \emph{(A)} End correction per end, in units of the radius, of a sharp-walled cylinder against $L/a$ (exact mode matching and its closed form): Hall's $\pi/4$ is the thin-plate limit; the exact value reaches the flanged-tube $0.8217$ by $L \approx a$. \emph{(B)} Deviation of three references from the converged map conductance at $L = 20$~nm: Equation~\eqref{eq:kowalczyk} at the area-conserving offset (the drift of Figure~\ref{fig:analytical_validation}), the exact access term without the rim correction, and Equation~\eqref{eq:semfull}. \emph{(C)} Error of Equation~\eqref{eq:kowalczyk} at the area-conserving offset over $d = 1.5$--$50$~nm and $L = 1$--$40$~nm. \emph{(D)} The rim length deficit, scaled by $t^{4/3}a_e^{-1/3}$, against $\lambda = (L+2t)/a_e$ for all $d = 2$--$50$~nm, and the corner-theory prediction $C\,[k(\lambda)/k_\infty]^{2}$ of Equation~\eqref{eq:rim}, which contains no fitted parameter.}
\label{fig:si_softwall}
\end{figure}

% ---------------------------------------------------------------------
% (3) Other sentences of the draft affected by this result
% ---------------------------------------------------------------------
% * Abstract: "brings SEM open-pore conductance within ~1% of a
%   parameter-free analytical reference over the design range d <= 20 nm"
%   -> "reproduces a parameter-free sharp-walled reference to 0.002%
%   (d >= 5 nm, L = 20 nm; 0.12% for d = 2-50 nm, L = 1-40 nm)".
% * Results 3 (Eq. kowalczyk): "it is accurate in the regime L >~ d" ->
%   Hall's access term is exact only as L -> 0; for L >~ d/2 it is 4.6%
%   small, and Eq. (kowalczyk) overestimates G by 0.8-3.2% for
%   L/d = 0.2-4. Keep Eq. (kowalczyk) as the familiar form, but benchmark
%   against Eq. (semfull).
% * Results 3.1: "The offset that would force point-by-point agreement ...
%   drifts from 0.52 nm at d = 5 nm to 0.95 nm at d = 25 nm as the Hall
%   access model leaves its validity regime" and "The two bounds ...
%   coincide at d ~ 20 nm, so we report quantitative agreement over
%   d <= 20 nm and treat 20-25 nm as a controlled extrapolation" -> the
%   drift is the access-term error (plus a small box term); no
%   extrapolation caveat is needed.
% * Appendix A (si_sigma caption/text): "this drift is not a failure of
%   the correction but the breakdown of the Hall thin-membrane access
%   model as d -> L" -> Hall's model is the thin-membrane limit; the drift
%   is its error for a thick membrane (4.6% in the access term), weighted
%   by the access share of the resistance, which grows with d.
% * Table "scope": rows "Solid-state geometries, d <= 20 nm" and
%   "d = 20-25 nm: Caution, controlled extrapolation (<= 3.5%)" -> one row,
%   d = 2-50 nm, L = 1-40 nm, "0.12% vs exact sharp-wall reference (0.002% at L = 20 nm, d >= 5 nm)".
% * Sec. domain / Appendix si_domain: the golden-aspect alpha = 1.2 was
%   chosen to "land on" Eq. (kowalczyk) at the corrected offset, which is
%   0.8% above the infinite-reservoir map conductance at d = 10 nm.
%   Against the exact reference the golden box is -0.24% to +0.54% from the
%   infinite-reservoir value (p/d = 0.35-2, d = 5-25 nm,
%   validation/soft_wall_mapping/golden_box.csv); \PENDING{re-run the
%   alpha sweep of Fig. si_box(C) against Eq. (semfull) and update the rule
%   if the optimum moves}.
% * Limitations: "the 0.54 nm offset is derived for the present
%   conductivity-ramp parameters and validated over d <= 25 nm at
%   L = 20 nm" -> the mapping holds for any linear ramp (r_cut = 3-6 A
%   tested) and d = 2-50 nm, L = 1-40 nm; tapered walls need the same
%   construction with their own exact reference (the axisymmetric solver
%   provides it for any piecewise-linear wall).
